% PAGES: 237-238
% Continues directly from sec07_c4_1.tex (PDF pages 235-236), which left the
% \begin{proof} of Lemma 7.13 open across the file boundary: the align* block
% ending "...= (C^{(1)})^t\Sigma_{\bfX}C^{(2)}," at the bottom of PDF page 236
% is completed in prose by the "since Sigma_X = ..." clause that opens PDF
% page 237 (folio 221) below -- closing that same sentence, not starting a
% new one. The proof environment is closed in this file, not reopened.

since $\Sigma_{\bfX} = \E{(\bfX-\E{\bfX})(\bfX-\E{\bfX})^t}$; see (7.131).

Since $\left(C^{(1)}\right)^t\Sigma_{\bfX}C^{(2)}$ is a number, it follows that
\[
\left(C^{(1)}\right)^t\Sigma_{\bfX}C^{(2)} \;=\;
\left(\left(C^{(1)}\right)^t\Sigma_{\bfX}C^{(2)}\right)^t;
\]

see also (7.28). Thus,
\[
\left(C^{(1)}\right)^t\Sigma_{\bfX}C^{(2)} \;=\; \left(C^{(2)}\right)^t\Sigma_{\bfX}^t
\left(\left(C^{(1)}\right)^t\right)^t \;=\; \left(C^{(2)}\right)^t\Sigma_{\bfX}C^{(1)},
\]
since $\Sigma_{\bfX}$ is a symmetric matrix and therefore $\Sigma_{\bfX}^t =
\Sigma_{\bfX}$.
\end{proof}

\begin{theorem}[Linear Transformation Property; see also Theorem 7.3.]
Let $\bfX$ be an $n \times 1$ multivariate random variable, and let $M$ be an $m \times
n$ matrix with constant entries. If $\bfY$ is the $m \times 1$ multivariate random
variable given by
\begin{equation}
\bfY \;=\; M\bfX,
\end{equation}
then
\begin{equation}
\Sigma_{\bfY} \;=\; M\Sigma_{\bfX}M^t,
\end{equation}
where $\Sigma_{\bfX}$ and $\Sigma_{\bfY}$ are the covariance matrices of $\bfX$ and of
$\bfY$, respectively.
\end{theorem}

\begin{proof}
From (7.131) and (7.137), it follows that
\begin{align}
\Sigma_{\bfY} &= \E{(\bfY-\E{\bfY})(\bfY-\E{\bfY})^t} \notag \\
&= \E{(M\bfX - \E{M\bfX})(M\bfX-\E{M\bfX})^t}.
\end{align}

From (7.128), we find that
\begin{equation}
M\bfX - \E{M\bfX} \;=\; M\bfX - M\E{\bfX} \;=\; M(\bfX-\E{\bfX}),
\end{equation}
and therefore
\begin{equation}
(M\bfX-\E{M\bfX})^t \;=\; (M(\bfX-\E{\bfX}))^t \;=\; (\bfX-\E{\bfX})^tM^t.
\end{equation}

From (7.139--7.141), and using (7.127), we obtain that
\begin{align*}
\Sigma_{\bfY} &= \E{M(\bfX-\E{\bfX})(\bfX-\E{\bfX})^tM^t} \\
&= M\,\E{(\bfX-\E{\bfX})(\bfX-\E{\bfX})^t}\,M^t \\
&= M\Sigma_{\bfX}M^t,
\end{align*}
since $\Sigma_{\bfX} = \E{(\bfX-\E{\bfX})(\bfX-\E{\bfX})^t}$; cf. (7.131).
\end{proof}

\section{References}

Details on the Linear Transformation Property and its applications to Monte Carlo
methods can be found in Andersen and Piterbarg \cite{3}.

An important practical issue is that, for large data sets, the covariance matrix of
time series data may be very close to singular, or some eigenvalues may be very small
and negative; see Lai and Xing \cite{25} for practical ways to overcome this issue.

Three different ways for finding all the values of a parameter $\rho$ such that a
matrix is a correlation matrix can be found in Stefanica, Radoi\v{c}i\'{c}, and Wang
\cite{39}.

For a comprehensive presentation of multivariate normal random variables, see Jacod and
Protter \cite{24}.

% PDF page 238 (folio 222) ends the References section here -- the last
% sentence ("...Jacod and Protter [24].") is complete, and the rest of the
% page is blank. No environment is left open. PDF page 239 (folio 223) opens
% fresh with the "7.9 Exercises" section heading at the top of the page,
% continued in sec07_c4_3.tex.
